\section{Finite-Codeword Collision Geometry}
\label{sec:collision-geometry}

Having defined relational fidelity independently of representation, we now study finite codewords as one concrete description family.
A hard code induces same-codeword equality and hence a partition into at most $K$ classes; probabilistic assignments relax this to the probability that independently drawn codewords collide.
This collision probability is the representation-side geometry used below, linking smooth optimization to hard partitions, structured binary parameterizations, and codeword-prevalence corrections.


\subsection{Probabilistic finite encoders}

A hard finite assignment is discrete, whereas learned representations are often optimized through continuous parameters.
We therefore begin by replacing each hard codeword with a categorical distribution over the same finite alphabet.
This retains the finite set of available states while providing a smooth representation of assignment uncertainty.

Let $\mathcal Z=\{1,\ldots,K\}$ and let an encoder map an input $x$ to
\begin{equation}
q_\theta(\cdot\mid x)\in\Delta^{K-1},\qquad
\Delta^{K-1}=\{q\in\R_{\ge0}^K:\textstyle\sum_z q_z=1\}.
\label{eq:discrete-encoder}
\end{equation}
The map $x\mapsto q_\theta(\cdot\mid x)$ may itself be deterministic.
Stochasticity enters only when a codeword is drawn from the resulting categorical distribution.

Consider inputs $x,x'$ associated with two carrier elements, and take conditionally independent draws
$Z\sim q_\theta(\cdot\mid x)$ and
$Z'\sim q_\theta(\cdot\mid x')$.
Their probability of receiving the same codeword is
\begin{equation}
k_\theta(x,x')
=\sum_zq_\theta(z\mid x)q_\theta(z\mid x')
=\Pr[Z=Z'\mid x,x'].
\label{eq:collision-kernel}
\end{equation}
Thus collision is a soft version of the equality relation supplied by a finite code.
It is large when the two assignment distributions place mass on the same codewords and becomes exact same-class membership when the distributions become deterministic.
Because this collision event involves two independently encoded inputs, it is the order-two, or pairwise, collision probability.
We write $k_{\theta,2}$ when the collision order needs to be explicit, but abbreviate it as $k_\theta$ throughout the main text.
More generally, an integer-order extension requires $\alpha\ge2$ independently encoded inputs to receive the same codeword; it is developed in \ref{app:higher-order-collision-geometry}.

The same expression also has a useful geometric interpretation.
It is the ordinary inner product of two categorical probability vectors, and is therefore a discrete probability-product kernel~\citep{jebara2004probability}.
In particular, it is symmetric, lies in $[0,1]$, and is positive semidefinite.
For any inputs $x_i$ and real coefficients $a_i$,
\begin{equation}
\sum_{i,j}a_i a_j k_\theta(x_i,x_j)
=\left\|\sum_i a_iq_\theta(\cdot\mid x_i)\right\|_2^2\ge0.
\end{equation}

Viewing collision as an inner product also shows what kinds of pairwise geometry this representation can express.
On a finite sample, let $Q$ contain the categorical assignment vectors as its rows, so the collision Gram matrix is $\Gamma_\theta=QQ^\top$.
Hence $\operatorname{rank}(\Gamma_\theta)\le K$: a $K$-codeword collision representation can express at most $K$ independent spectral directions of the pairwise relation.
Raw collision therefore represents a restricted family of relational structures.
The broader framework also includes decoders such as the block-valued construction in \cref{sec:block-summary-walkthrough}, which need not be limited to this particular form.

There is one further distinction specific to probabilistic assignments.
The diagonal of the collision kernel is
\begin{equation}
k_\theta(x,x)=\sum_zq_\theta(z\mid x)^2
=\|q_\theta(\cdot\mid x)\|_2^2\in[1/K,1].
\label{eq:self-collision}
\end{equation}
For a fixed input $x$, let $Z_1,Z_2\overset{\mathrm{iid}}{\sim}q_\theta(\cdot\mid x)$.
Then $k_\theta(x,x)=\Pr[Z_1=Z_2\mid x]$, so the diagonal of the collision kernel measures the concentration of the conditional assignment distribution.
A uniform assignment gives $k_\theta(x,x)=1/K$, while a deterministic assignment gives $k_\theta(x,x)=1$.
Hence the diagonal entries of a soft collision matrix lie in $[1/K,1]$, and become one in the deterministic hard-partition limit.


\subsection{Hard partitions and their relaxation}

The probabilistic formulation is useful for optimization, but the finite representation of primary interest is often a hard code.
The connection between the two is direct.
A deterministic assignment $c_\theta:\mathcal X\to\mathcal Z$ is represented by the one-hot distribution
$q_\theta(z\mid x)=\ind[z=c_\theta(x)]$.
Substituting this distribution into \cref{eq:collision-kernel} gives
\begin{equation}
k_\theta(x,x')=\ind[c_\theta(x)=c_\theta(x')].
\label{eq:hard-collision}
\end{equation}
The collision relation is then exactly membership in the same part of the partition induced by $c_\theta$.

The same conclusion holds as a limit.
Let $\{q_n(\cdot\mid x)\}_{n\ge1}$ be a sequence of categorical assignment distributions that, for every input under consideration, converges to the one-hot distribution at $c_\theta(x)$.
Then
\[
\sum_z q_n(z\mid x)q_n(z\mid x')
\longrightarrow
\ind[c_\theta(x)=c_\theta(x')].
\]
Collision can therefore be viewed as a continuous relaxation of hard same-partition membership.

At nonzero assignment uncertainty, collision is a graded affinity and generally lacks the transitivity of an equivalence relation.
The distinction between the soft assignment geometry and the hard representation is consequential: a soft encoder can distribute mass broadly across the alphabet even when a chosen hardening rule assigns many or all inputs to the same codeword.
Soft and hard relational statistics should therefore be interpreted separately.


\subsection{Factorized binary codes}

A categorical alphabet with $K$ states becomes expensive to represent explicitly when $K$ is large.
A useful specialization obtains an exponentially large finite alphabet from a small number of binary coordinates.
Let the code contain $b$ conditionally independent bits, so that
$\mathcal Z=\{0,1\}^b$ and $K=2^b$.
Writing
$p_r(x)=\Pr[Z_r=1\mid x]$ for the conditional probability that bit $r$ takes value one at input $x$, the probability of a complete binary word
$a=(a_1,\ldots,a_b)$ is
\begin{equation}
q_\theta(a\mid x)
=\prod_{r=1}^bp_r(x)^{a_r}(1-p_r(x))^{1-a_r},
\qquad a\in\{0,1\}^b.
\label{eq:factorized-code}
\end{equation}

The main computational benefit is that collision can be evaluated without enumerating all $2^b$ complete words.
Summing the equality probability over the factorized code distribution gives
\begin{equation}
k_\theta(x,x')
=\prod_{r=1}^b\left[p_r(x)p_r(x')+(1-p_r(x))(1-p_r(x'))\right].
\label{eq:bit-collision}
\end{equation}
Each factor is the probability that the two independently sampled values of bit $r$ agree.
Their product is therefore the probability that the \emph{entire} binary code agrees.
As every $p_r(x)$ approaches zero or one, the expression reduces to deterministic equality of the complete codewords.

The factorization assumes that the bits are independent conditional on each fixed input $x$, which parameterizes $q_\theta(\cdot\mid x)$ as a product distribution.
Across the input population, the probabilities $p_r(x)$ may co-vary, so the aggregate codeword distribution need not factorize across bit coordinates.

The binary parameterization provides a compact factorization of a categorical distribution over $K=2^b$ complete codewords.
Collision treats each complete word as one categorical state, so its categorical representation lives in $K$ dimensions.
The spherical construction developed later applies to these categorical probability vectors.


\subsection{Reference-weighted codeword affinities}

Raw collision treats agreement on every codeword equally.
Some relational objectives instead require agreement to be weighted by codeword prevalence.
Agreement on a codeword used by almost every element may carry less distinguishing information than agreement on a rarely occupied word.
This motivates measuring a shared assignment relative to a reference prevalence for that codeword.

Let $r=(r(z))_{z\in\mathcal Z}$ be a fixed positive reference distribution on $\mathcal Z$.
Define the reference-weighted affinity
\begin{equation}
a_r(x,x')
=\sum_z\frac{q_\theta(z\mid x)q_\theta(z\mid x')}{r(z)}
=\left\langle\frac{q_\theta(\cdot\mid x)}{\sqrt r},
\frac{q_\theta(\cdot\mid x')}{\sqrt r}\right\rangle.
\label{eq:reference-affinity}
\end{equation}
The square root and division in the vector expression are taken coordinatewise.
Like raw collision, this quantity is positive semidefinite because it is an inner product after coordinate rescaling.
Inverse-reference weighting can make this value exceed one, so we treat it as an affinity rather than a probability.

The hard-code interpretation makes the weighting especially transparent.
If two inputs receive the same deterministic word $z$, their affinity is $1/r(z)$; if they receive different words, it is zero.
Agreement on a word with small reference mass is therefore weighted more strongly than agreement on a common word.
For the uniform reference $U_K(z)=1/K$, every word receives the same correction and
$a_{U_K}=Kk_\theta$, so the weighting changes only the global scale.

A particularly important reference is supplied by the encoder's own aggregate codeword distribution.
Fix a probability measure $\mu$ on carrier-element inputs and define
\begin{equation}
q_Z(z)
:=\Pr(Z=z)
=\E_{X\sim\mu}[q_\theta(z\mid X)].
\label{eq:aggregate-code}
\end{equation}
This distribution records how frequently each codeword is used on average under the declared population.
Taking the reference to be this aggregate gives
\begin{equation}
a_{q_Z}(x,x')
=\sum_{z:q_Z(z)>0}\frac{q_\theta(z\mid x)q_\theta(z\mid x')}{q_Z(z)}.
\label{eq:qz-affinity}
\end{equation}
The affinity now corrects agreement by the actual prevalence of the shared codeword: collisions through frequently occupied states are downweighted, while collisions through less common states receive greater weight.

Inactive codewords are omitted from the sum.
Indeed, $q_Z(z)=0$ implies
$q_\theta(z\mid X)=0$ almost surely under $\mu$, so this support convention is well defined on the declared population.
Inputs outside the support of that population would require a separate convention.

This endogenous inverse-mass correction will be important later.
The same form appears in the exact centroid-reconstruction and normalized-association correspondences, where codeword mass determines how within-code agreement should be normalized.
An \emph{external} reference can separately specify a desired marginal organization when used in a prior-matching criterion.

\section{Codeword Organization and Spherical Geometry}
\label{sec:collision-entropy}

Collision also describes how a finite code is organized across a population: averaging it yields aggregate occupancy and R\'enyi-2 effective code count.
We then distinguish external-reference organization from endogenous mass correction and use radial normalization to expose the spherical geometry of categorical assignments.

\subsection{Population collision and effective occupancy}

For $X,X'\overset{\mathrm{iid}}{\sim}\mu$, the population collision probability depends only on the aggregate codeword distribution:
\begin{align}
\E[k_\theta(X,X')]
&=\sum_z\E[q_\theta(z\mid X)]\E[q_\theta(z\mid X')]\\
&=\sum_zq_Z(z)^2.
\label{eq:population-collision}
\end{align}
The right-hand side is the collision probability of the aggregate codeword distribution itself.
It is large when population mass is concentrated on a few codewords and smaller when usage is spread more broadly.
The iid population sampling is essential to this occupancy interpretation: independence makes the average collision depend only on the aggregate codeword probabilities $q_Z(z)$.
For pairs drawn from a relation-specific joint distribution, such as graph edges, the average collision also reflects dependence between the paired assignments and is therefore a different relational statistic.
The measure $\mu$ must therefore be part of the interpretation: it may be uniform over the carrier of a single object, degree-weighted over the vertices of a graph, or defined by a declared pooling of carrier elements across a collection of relational objects.

Because population collision is exactly the squared $\ell_2$ mass of the aggregate codeword distribution, it has a natural information-theoretic representation through order-two R\'enyi entropy~\citep{renyi1961measures}:
\begin{equation}
H_2(q_Z)=-\log\sum_zq_Z(z)^2
=-\log\E[k_\theta(X,X')].
\label{eq:collision-entropy}
\end{equation}
We write $R_2=H_2(q_Z)$ for this collision-occupancy measure.
Entropy is useful analytically, but the same quantity can be expressed as an effective number of equally occupied codewords:
\begin{equation}
K_{\mathrm{eff}}=\exp(H_2(q_Z))
=\frac{1}{\sum_zq_Z(z)^2},\qquad 1\le K_{\mathrm{eff}}\le K.
\label{eq:effective-code-count}
\end{equation}
Uniform use of exactly $M$ words gives $K_{\mathrm{eff}}=M$.
For nonuniform occupancy, the same expression gives the inverse-collision, or inverse-Simpson, effective count, which decreases as probability mass becomes more concentrated.

Collision-effective occupancy is only one summary of code use.
It is useful to place it alongside nominal alphabet size, active support, and Shannon entropy because these quantities coincide only in special cases:
\begin{equation}
\begin{aligned}
R_{\mathrm{nom}}&=\log K,&
R_0&=\log|\operatorname{supp}(q_Z)|,\\
R_1&=-\sum_zq_Z(z)\log q_Z(z),&
R_2&=-\log\sum_zq_Z(z)^2.
\end{aligned}
\label{eq:finite-state-rate-hierarchy}
\end{equation}
They satisfy $R_2\le R_1\le R_0\le R_{\mathrm{nom}}$.
All logarithms are natural unless specified otherwise, and dividing these quantities by $\log 2$ expresses them in bits.
These quantities summarize different aspects of codeword use.
When codewords are encoded under an appropriate probabilistic coding model, the Shannon entropy $H_1(q_Z)$ can also characterize the asymptotic average coding cost per codeword.
Any such encoded payload that forms part of the retained description contributes to the resource functional $\Omega$.

Aggregate occupancy also differs from conditional assignment uncertainty.
If every input has a deterministic assignment, self-collision is one even when population collision is small.
Conversely, the input-independent encoder $q_\theta(\cdot\mid x)=U_K$ has maximal aggregate $H_2$ but conveys no input dependence through its sampled codeword.
For soft encoders, broad aggregate occupancy can therefore coexist with a concentrated hard partition or weak input dependence.
The associated conditional-collision and divergence-based dependence quantities are recorded in \ref{app:conditional-collision-dependence}.

\subsection{External references and endogenous mass correction}

Aggregate collision measures concentration without specifying what pattern of codeword use is preferred.
In some settings, organization is meaningful only relative to a reference distribution.
An external reference can encode a desired marginal pattern, such as equal use of exchangeable codewords, whereas choosing the reference masses as $r(z)=q_Z(z)$ serves a different purpose: it compensates local agreement for the prevalence actually produced by the encoder.
For fixed positive reference masses $r(z)$, independent population averaging gives
\begin{equation}
\E[a_r(X,X')]=\sum_z\frac{q_Z(z)^2}{r(z)},
\qquad
D_2(q_Z\|r)=\log\E[a_r(X,X')].
\label{eq:renyi-reference-affinity}
\end{equation}
Thus reference-weighted affinities have a population interpretation through R\'enyi divergence~\citep{renyi1961measures,vanerven2014renyi}.
When the corresponding divergence is used as an organization penalty, the external reference specifies the target marginal organization of the alphabet.
For exchangeable words, uniform matching gives
\begin{equation}
D_2(q_Z\|U_K)=\log K-H_2(q_Z).
\label{eq:uniform-renyi}
\end{equation}
Minimizing this divergence encourages broader effective occupancy.
Adding $+\lambda H_2(q_Z)$ to a minimized objective does the opposite: it encourages concentration.
The sign of the regularizer therefore determines whether it encourages broad occupancy or concentration.

For comparison, $D_{\mathrm{KL}}(q_Z\|U_K)=\log K-H_1(q_Z)$ also has the uniform optimum.
Although $H_2\le H_1$ for each distribution, the two entropies can rank candidate distributions differently and induce distinct level sets around their common optimum.
The general identity $D_\alpha(q\|U_K)=\log K-H_\alpha(q)$ and the sign conventions are collected with the higher-order material in \ref{app:higher-order-collision-geometry}.

For the endogenous choice $r(z)=q_Z(z)$, inverse-mass normalization gives
\begin{equation}
\E[a_{q_Z}(X,X')]=1.
\label{eq:endogenous-reference-expectation}
\end{equation}
For almost every fixed $x$, even $\E_{X'}[a_{q_Z}(x,X')]=1$.
This unit expectation results from correcting local agreement by the occupancy induced by the encoder.
The aggregate masses $q_Z(z)$ therefore supply endogenous inverse-mass normalization when used as the reference and can separately enter a prior-matching penalty when compared with an external reference.

\subsection{Radial normalization and cosine similarity}

The raw collision kernel is the Euclidean inner product of categorical probability vectors, and therefore combines two geometric effects: the angular agreement between their directions and the magnitudes of their Euclidean norms.
For categorical distributions, these norms are themselves meaningful, since their squared values are self-collision probabilities and hence encode R\'enyi-2 concentration.
It is therefore useful to separate concentration from angular agreement when studying the geometry induced by collision.
Radial normalization does exactly this, mapping categorical probability vectors to the positive unit sphere, where normalized collision becomes cosine similarity.
For $q\in\Delta^{K-1}$, define
\begin{equation}
u(q)=\frac{q}{\|q\|_2}\in\mathbb S_+^{K-1},\qquad
\mathbb S_+^{K-1}=\{u\in\R_{\ge0}^K:\|u\|_2=1\}.
\label{eq:sphere-map}
\end{equation}
The map is bijective, with inverse
\begin{equation}
q=\frac{u}{\|u\|_1}.
\label{eq:sphere-map-inverse}
\end{equation}
A more standard spherical representation of probability distributions is the elementwise square-root map $q\mapsto\sqrt q$, for which the spherical inner product becomes the Bhattacharyya coefficient and which underlies Hellinger and information-geometric constructions~\citep{amari2016information}.
The radial normalization used here serves a different purpose: because collision is the ordinary inner product $q^\top q'$, normalizing by the Euclidean norm isolates the angular component of that collision geometry while preserving a bijective representation of the simplex.

Writing $k_2(q,q')=q^\top q'$, normalized collision is
\begin{equation}
u(q)^\top u(q')
=\frac{k_2(q,q')}{\sqrt{k_2(q,q)k_2(q',q')}}.
\label{eq:normalized-collision}
\end{equation}
Thus normalized collision is exactly cosine similarity on the nonnegative spherical orthant: it keeps the angular agreement between assignments, while raw collision also depends on the concentration of each assignment.\footnote{The negative logarithm of this normalized inner product is the discrete analogue of the Cauchy--Schwarz divergence of \citet{jenssen2006cauchy}.}
This connects the finite-codeword realization to hyperspherical alignment~\citep{wang2020alignment}, while retaining the categorical constraints.
The one-hot representations of hard words occupy the fixed coordinate directions $e_z$.
Relabeling complete codewords permutes these directions; it does not create an arbitrary rotationally invariant embedding family.
A factorized bit parameterization can impose further restrictions on which codeword relabelings preserve its parameterized family.

The normalized similarity factors out the explicit self-collision multipliers from the pairwise score.
\begin{equation}
\|q\|_2^2
=k_2(q,q)
=\sum_z q(z)^2
=\exp[-H_2(q)],
\qquad
\|q\|_2
=\exp[-H_2(q)/2].
\label{eq:collision-norm-entropy}
\end{equation}
If $\cos\vartheta=u(q)^\top u(q')$, then
$k_2(q,q')=q^\top q'=\|q\|_2\|q'\|_2\cos\vartheta$, and therefore
\begin{equation}
k_2(q,q')
=\exp\!\left[-\tfrac12\bigl(H_2(q)+H_2(q')\bigr)\right]\cos\vartheta.
\label{eq:collision-angular-decomposition}
\end{equation}
This decomposition shows that raw collision combines angular agreement with the R\'enyi-2 concentration of each assignment.
The spherical map itself is lossless on the simplex, since $q$ can be recovered from $u(q)$ through \cref{eq:sphere-map-inverse}, while cosine similarity isolates only the angular component.
Thus raw and normalized collision emphasize different aspects of the same categorical geometry.

\subsection{Entropy in spherical coordinates}

Having separated raw collision into angular agreement and R\'enyi-2 concentration in \cref{eq:collision-angular-decomposition}, we now express the entropy-based occupancy measures introduced in \cref{eq:finite-state-rate-hierarchy} in the same spherical coordinates.
This provides a common geometric language for two aspects of the finite-codeword representation: angles between spherical vectors describe normalized collision between assignments, while the geometry of an individual spherical vector describes the concentration and organization of the corresponding categorical distribution.
For $u=q/\|q\|_2$, the inverse relation gives
\begin{equation}
H_2(q)=2\log\|u\|_1.
\label{eq:h2-spherical}
\end{equation}
Let $u_U=\mathbf1/\sqrt K$ and let $\vartheta_U$ be its angle to $u$.
Since $\cos\vartheta_U=\|u\|_1/\sqrt K$,
\begin{equation}
D_2(q\|U_K)=-2\log\cos\vartheta_U
=\log\bigl(1+K\|q-U_K\|_2^2\bigr).
\label{eq:renyi-angular}
\end{equation}
These are exact changes of coordinates.
At the level of one input, increasing $H_2(q_\theta(\cdot\mid x))$ makes the assignment less concentrated.
At the level of the aggregate probabilities $q_Z(z)$, increasing $H_2(q_Z)$ spreads aggregate codeword occupancy.
The same algebra describes different objects depending on which distribution is inserted.

R\'enyi-2 has an especially simple angular description, while Shannon entropy provides another organization measure expressible through $u$:
\begin{equation}
H_1(q)=\log\|u\|_1
-\frac{1}{\|u\|_1}\sum_i u_i\log u_i,
\label{eq:shannon-spherical}
\end{equation}
with $0\log0=0$.
Unlike $H_2$, it depends on the full coordinate profile rather than only the angle to the uniform direction.
The two entropies share collapsed and uniform extrema but generally have different level sets.
The broader R\'enyi expression appears in \ref{app:renyi-spherical-coordinates}.

Taken together, these identities explain why order two provides a particularly coherent description of this finite-codeword realization: local equality probability, population collision, and normalized cosine geometry all arise from the same pairwise collision statistic.
The empirical and modelling choice between R\'enyi-2, Shannon, and other regularizers remains application dependent.

\section{Alignment, Separation, and Relational Requirements}
\label{sec:alignment-separation}

We now use the collision geometry to express source-specific relational requirements.
Favored and disfavored pairs control alignment and separation, while an optional marginal term controls how the alphabet is used; these distinct roles lead to the learning objectives below.

\subsection{From the source relation to a learning objective}

Let $(\mathsf V,\mathsf V')\sim P_{\mathrm{pair}}$ be a declared distribution on carrier pairs, and let $x_{\mathsf V},x_{\mathsf V'}$ denote their corresponding encoder-side observations.
A code-induced pair statistic $\psi_\theta$ may be raw collision, reference-weighted affinity, normalized collision, or a dissimilarity such as $-\log k_\theta$.
A generic pairwise learning objective is
\begin{equation}
\mathcal L_{\mathrm{rel}}(\theta)
=\E_{(\mathsf V,\mathsf V')\sim P_{\mathrm{pair}}}\left[\ell\bigl(S(\mathsf V,\mathsf V'),\psi_\theta(x_{\mathsf V},x_{\mathsf V'})\bigr)\right].
\label{eq:generic-relational-objective}
\end{equation}
The statistic may itself be the reconstructed relation, or it may control a decoder or selector used to form that relation.
Whether this objective equals the declared relational distortion depends on that reconstruction or evaluation convention; otherwise it may serve as a training surrogate for the desired fidelity.

For a favored-pair distribution $P_+$ on carrier pairs, an alignment loss decreases with the desired affinity:
\begin{equation}
\mathcal L_{\mathrm{align}}
=\E_{(\mathsf V,\mathsf V')\sim P_+}[\ell_+(\psi_\theta(x_{\mathsf V},x_{\mathsf V'}))].
\label{eq:general-alignment}
\end{equation}
For example, $\ell_+(k)=-\log k$ promotes collision, with infinite loss at zero overlap unless a numerical surrogate is explicitly substituted.
Using normalized collision instead promotes angular agreement without the same concentration factors.
For dissimilarity-valued statistics, the monotonic directions reverse.

\subsection{Pair-specific separation and aggregate organization}

A source can also identify pairs whose collision would erase a required distinction.
For an appropriate disfavored-pair distribution $P_-$ on carrier pairs,
\begin{equation}
\mathcal L_{\mathrm{sep}}^{\mathrm{rel}}
=\E_{(\mathsf V,\mathsf V')\sim P_-}[\ell_-(\psi_\theta(x_{\mathsf V},x_{\mathsf V'}))],
\label{eq:relational-separation}
\end{equation}
where $\ell_-$ increases with undesirable affinity.
This is a pair-specific condition.
By contrast,
\begin{equation}
\mathcal L_{\mathrm{org}}=D_2(q_Z\|r)
\label{eq:marginal-organization}
\end{equation}
compares population occupancy with an external reference.
Two partitions can have identical occupancy and different pairwise equivalences: marginal matching organizes population use, while pair-specific terms specify which source relations should survive.

A combined objective can therefore be written as
\begin{equation}
\mathcal L
=\mathcal L_{\mathrm{align}}
+\lambda_{\mathrm{rel}}\mathcal L_{\mathrm{sep}}^{\mathrm{rel}}
+\lambda_{\mathrm{org}}D_2(q_Z\|r),
\qquad\lambda_{\mathrm{rel}},\lambda_{\mathrm{org}}\ge0.
\label{eq:alignment-relational-marginal}
\end{equation}
Only the terms required by the learning problem need be present.
Uniform prior matching supplies one possible population-level noncollapse mechanism, while relational separation acts on specified pairs.
This clarifies the analogy with alignment and uniformity on a hypersphere~\citep{wang2020alignment}: the finite categorical case has a fixed available frame and an explicitly chosen marginal reference, rather than an unrestricted continuous geometry.

\subsection{Occupancy, fidelity, and resource constraints}
\label{sec:occupancy-frontiers}

Occupancy is a property of realized code use, not by itself the complete description resource.
We therefore study occupancy constraints alongside the relational distortion defined in \cref{sec:relational-distortion} within the chosen finite-codeword family.
Let $m_{S,\theta}$ denote the description induced by parameter setting $\theta$, and define $D_{\mathrm{rel}}(\theta):=\Delta_V\!\left(S,\mathsf{Dec}_V(m_{S,\theta})\right)$ as the relational distortion within this family, following the complete-description formulation in \cref{eq:lossy-relational-compression}.
Then
\begin{equation}
D_{\mathrm{rel}}^\star(K)=\inf_{\theta:|\mathcal Z|\le K}D_{\mathrm{rel}}(\theta).
\end{equation}
The other parts of the description and decoder family must be held fixed or explicitly controlled; otherwise an auxiliary payload could carry arbitrary source information while the code alphabet remains small.

For $R_2(\theta)=H_2(q_Z)$, define
\begin{equation}
D_{\le}^\star(R)=\inf_{\theta:R_2(\theta)\le R}D_{\mathrm{rel}}(\theta),\qquad
D_{\ge}^\star(R)=\inf_{\theta:R_2(\theta)\ge R}D_{\mathrm{rel}}(\theta).
\label{eq:relational-complexity-distortion-frontier}
\end{equation}
An upper occupancy bound and a lower occupancy requirement answer different questions.
The first limits realized code use; the second enforces noncollapse or organization.
An inverse requirement can also be written $R_2^\star(D)=\inf_{\theta:D_{\mathrm{rel}}(\theta)\le D}R_2(\theta)$.
These definitions are occupancy--fidelity relations; an operational coding model can additionally relate them to code length.

For direct positive-relation distortion, the one-codeword solution can have zero distortion, making the upper-bound problem trivial.
A lower bound on $R_2$ instead asks how much relational fidelity can be retained while using a noncollapsed alphabet.
At fixed $K$, an objective
\begin{equation}
D_{\mathrm{rel}}(\theta)+\lambda_{\mathrm{org}}D_2(q_Z\|U_K)
=D_{\mathrm{rel}}(\theta)-\lambda_{\mathrm{org}}R_2(\theta)
+\lambda_{\mathrm{org}}\log K
\label{eq:organization-lagrangian}
\end{equation}
encourages this second direction.
Because this objective encourages occupancy rather than imposing an upper resource budget, its empirical distortion trajectory can increase with $R_2$.
Nonconvex optimization and discrete feasible sets also mean that scalar penalties need not recover every constrained frontier point.
